\chapter{Canales, salida al mercado y ventas}\label{ch:channels-gtm-and-sales}

% Alcance: estrategia de canales (la 4.ª P), movimiento GTM PLG frente a liderado por ventas,
%   dimensionamiento y compensación de la fuerza de ventas, ventas B2B y gestión del pipeline.
%   Capítulo nuevo (cierra la brecha de canales + ventas/GTM).
% Se sitúa aguas arriba del funnel digital (\cref{ch:digital-funnel-and-crm}): dónde se vende
%   determina qué funnel se construye.
% Vínculos cruzados: \cref{eq:cac}, \cref{eq:ltv-cac}, \cref{eq:viral-k}, \cref{eq:pipeline-ev},
%   \cref{eq:littles-law}, \cref{sec:kerr-trap}, \cref{eq:segment-score}.

\section{Estrategia de canales y la 4.ª P}\label{sec:channel-strategy}
% complejidad: 4

¿Cómo debe una empresa hacer su producto física y comercialmente disponible para los clientes?
La distribución --- la cuarta «P» de la mezcla de \emph{marketing} --- determina quién asume el
costo de llegar al cliente, quién es dueño de la relación y, en última instancia, cuánto del
margen unitario retiene la empresa. La decisión suele tratarse como una cuestión logística, pero
es, antes que nada, una cuestión de margen: cada intermediario añadido a la cadena reduce el
precio que recibe el fabricante y, a cambio, transfiere una porción del costo de alcance,
inventario y servicio. La pregunta estratégica es si el margen que se cede es menor que el
costo fijo de llegar a esos clientes de forma directa.

Sea \(P\) el precio para el usuario final, \(C_{\text{var}}\) el costo variable por unidad, \(Q\) el
volumen vendido y \(\Phi_{\text{fixed}}\) el costo fijo del canal directo (un equipo de ventas,
una organización de soporte, infraestructura de transacciones). La contribución de cada ruta es:
\begin{equation}\label{eq:channel-margin}
  \Pi_{\text{direct}}   = (P - C_{\text{var}})\,Q, \qquad
  \Pi_{\text{indirect}} = (P_w - C_{\text{var}})\,Q,
\end{equation}
donde \(P_w < P\) es el precio mayorista pagado por el intermediario. Tras descontar el costo fijo
del canal directo, la empresa prefiere la distribución directa cuando
\(\Pi_{\text{direct}} - \Phi_{\text{fixed}} > \Pi_{\text{indirect}}\), o de forma equivalente cuando
\((P - P_w)\,Q > \Phi_{\text{fixed}}\): la prima de margen unitario de ir directo, escalada por
el volumen, debe superar el costo fijo de operar el canal directo. A bajo volumen gana el
intermediario; a alto volumen gana la ruta directa. El punto de cruce es el volumen de umbral de
rentabilidad del canal.

\begin{figure}[htbp]
  \centering
  \includegraphics[width=0.78\linewidth]{channel-structure}
  \caption{Tres estructuras de canal: directa (fabricante a cliente), de un nivel (a través de un
    único revendedor o VAR) y de dos niveles (a través de un distribuidor más revendedor). Cada
    capa de intermediario captura un diferencial de margen conforme a \cref{eq:channel-margin},
    reduciendo el ingreso neto unitario del proveedor. La estructura óptima minimiza el costo
    total de distribución con la cobertura de mercado requerida.}
  \label{fig:channel-structure}
\end{figure}

\begin{example}[Decisión entre canal directo y revendedor]
La empresa SaaS de referencia cobra \money{50}/mes (\money{600}/año). Su costo variable es
\money{29}/mes (alojamiento \money{15} más soporte, procesamiento de pagos y éxito del cliente
\money{14}), lo que deja un margen de contribución de \money{21}/mes en el canal directo.

Un revendedor regional propone tomar una participación de margen del \qty{30}{\percent} sobre el
precio al usuario final, de modo que el proveedor recibe un precio mayorista de
\(P_w = \money{50}\times0.70 = \money{35}/\)mes. El margen de contribución indirecto por cliente
es entonces \(\money{35} - \money{29} = \money{6}/\)mes --- un tercio del margen directo. El
sacrificio de margen es \(\money{21} - \money{6} = \money{15}\)/mes por cliente.

Suponga que la empresa debe gastar \money{90000}/mes en costo fijo de ventas directas y de éxito
del cliente para atender el territorio del revendedor sin él. El número incremental de clientes en
el que el canal directo resulta preferible es:
\[
  Q^* \;=\; \frac{\Phi_{\text{fixed}}}{P - P_w} \;=\; \frac{\money{90000}}{\money{15}}
       \;=\; \num{6000}\ \text{clientes}.
\]
La empresa atiende actualmente alrededor de \num{6667} clientes de pago (\money{4000000} ARR
\(\div\) \money{600}/año). A esa escala, el canal directo aporta
\(\num{6667}\times\money{21} = \money{140000}\)/mes tras los costos variables, frente a
\(\num{6667}\times\money{6} = \money{40000}\)/mes de forma indirecta; la brecha de \money{90000}
en costo fijo queda cubierta. La implicación: el acuerdo con el revendedor tenía sentido económico
cuando la base instalada estaba por debajo de \num{6000}, pero la empresa ya ha cruzado el umbral
donde construir un canal directo es rentable. Si el revendedor abre una nueva geografía que la
empresa no puede alcanzar a costo comparable, el canal indirecto sigue creando valor allí --- el
cálculo del umbral de rentabilidad se aplica segmento por segmento, no a nivel de compañía.
\end{example}

\begin{admonition}{Advertencia}
Operar simultáneamente un equipo de ventas directo y un revendedor en el mismo territorio crea
\textit{conflicto de canal de distribución dual}: el revendedor compite por el mismo prospecto,
hace descuentos para ganar y, con el tiempo, enseña a los clientes a esperar el precio indirecto
más bajo. La empresa se ve entonces ante la disyuntiva de proteger la relación con el socio de
canal o defender el precio directo. Los acuerdos de canal que definen exclusividad territorial o
derechos de registro de \emph{deal} reducen este riesgo, pero sólo si se aplican de forma
consistente; la aplicación inconsistente es peor que no tener ninguna política, porque genera
incertidumbre para todas las partes.
\end{admonition}

\cref{eq:channel-margin} aplica igualmente a bienes físicos y a software, pero la estructura de
costos difiere marcadamente: el \(C_{\text{var}}\) de una empresa SaaS está dominado por el
alojamiento y la mano de obra de éxito del cliente, no por el inventario, de modo que el umbral
de costo fijo \(\Phi_{\text{fixed}}\) --- el costo de una organización de ventas directas ---
domina el cálculo a bajo volumen. Esa asimetría es lo que hace que el movimiento liderado por
producto examinado en \cref{sec:gtm-motion} sea tan poderoso en términos económicos: construye
un canal directo con un costo fijo casi nulo por cliente adquirido. \textcite{coughlan2006marketing}
desarrollan el marco completo de diseño de canales; \textcite{kotler2021} lo sitúa dentro de la
mezcla de \emph{marketing} integrada.


\section{PLG frente a movimiento GTM liderado por ventas}\label{sec:gtm-motion}
% complejidad: 4

Una vez establecida la estructura de canales, la empresa debe decidir cómo entra la demanda en el
canal: ¿el producto se vende a sí mismo mediante pruebas gratuitas y una incorporación
(\emph{onboarding}) de autoservicio, o un equipo de ventas conduce todo el recorrido desde la
conciencia hasta el cierre? Esta es la decisión sobre el movimiento de salida al mercado, y es una
de las decisiones estratégicas de mayor impacto que toma una empresa de software en etapa temprana,
porque determina la forma de \cref{eq:cac} --- el denominador de \cref{eq:ltv-cac} --- antes de
que la empresa haya gastado un dólar en adquisición.

En un movimiento de \textit{crecimiento liderado por producto} (PLG), el propio producto es el
principal vehículo de adquisición: los usuarios se registran, experimentan el valor central dentro
de una prueba de autoservicio y se convierten en clientes de pago sin intervención humana. El CAC
es bajo porque no hay compensación de ventas incorporada en el costo de adquisición; el coeficiente
viral (\cref{eq:viral-k}) puede multiplicar la base instalada a un costo marginal casi nulo cuando
los usuarios existentes invitan a colegas. PLG domina cuando el producto entrega valor inmediato y
tangible a un usuario individual que puede tomar o influir en una decisión de compra de forma
autónoma --- típicamente compradores de pequeñas y medianas empresas (PYME) con valores de contrato
anual (ACV) inferiores a \money{5000}.

En un movimiento \textit{liderado por ventas}, el producto está restringido: un prospecto debe
interactuar con un representante de ventas antes de acceder a una demostración significativa. El
CAC es elevado porque incluye el salario base, la comisión y el costo indirecto de un SDR que
califica el \emph{lead} antes de transferirlo a un ejecutivo de cuenta. El movimiento liderado
por ventas se justifica cuando la decisión involucra a múltiples partes interesadas, una revisión
de seguridad o legal, o un proceso de adquisición que una compra de autoservicio no puede
gestionar --- típicamente ACV empresarial superior a \money{50000}. \textcite{ross2011predictable}, la referencia profesional canónica para las ventas \emph{outbound} en SaaS, formalizó la separación SDR/AE/CSM que este capítulo asume como estándar.

El movimiento híbrido, \textit{ventas lideradas por producto} (PLS), superpone un equipo de
asistencia en ventas sobre un funnel de autoservicio. El catalizador de productividad es el
\textit{lead} \textit{calificado por el producto} (PQL): un usuario activo que ha cruzado un
umbral de participación demostrable se convierte en la señal que activa la prospección del SDR.
Dado que el representante de ventas contacta a un usuario que ya ha experimentado el valor, la
brecha de confianza se cierra antes de la primera conversación. La lógica de segmentación de
\cref{eq:segment-score} se aplica directamente a la elección del movimiento: los segmentos con
alto potencial de LTV pero baja conversión de autoservicio son candidatos a PLS; los segmentos con
alta conversión de autoservicio y bajo LTV son candidatos a PLG puro.

\begin{figure}[htbp]
  \centering
  \includegraphics[width=0.86\linewidth]{plg-vs-slg}
  \caption{El espectro PLG--SLG. Al desplazarse hacia la derecha, el CAC sube, el ACV sube y el
    ciclo de ventas se alarga. El disparador del PQL (línea discontinua) marca la transición del
    autoservicio a la asistencia en ventas; suele ser un hito de participación dentro del
    producto, no una puntuación de calificación de \emph{marketing}.}
  \label{fig:plg-vs-slg}
\end{figure}

\begin{example}[Modelo híbrido PLG + SDR]
La empresa adquiere actualmente \num{200} clientes por mes a través de un funnel de autoservicio
con \(\text{CAC} = \money{600}\) (de \cref{eq:cac} en \cref{ch:customer-economics}).
Segmentar la base de usuarios revela dos cohortes distintas: cuentas de PYME que se incorporan en
autoservicio con CAC\,\(\approx\)\,\money{300}, y cuentas de mercado medio que se estancan antes
de convertir. Un movimiento de SDR dirigido a los PQL de mercado medio tiene un
CAC\,\(\approx\)\,\money{800}, pero el LTV de mercado medio es aproximadamente \(3\times\) el de
PYME, de modo que la razón de \cref{eq:ltv-cac} aún supera el umbral de \(3\times\).

El disparador de PQL elegido es \textit{más de cinco informes generados en los primeros catorce
días de la prueba}. Los usuarios que alcanzan este umbral convierten al \qty{45}{\percent}; la
tasa de conversión de referencia para todos los usuarios de prueba es del \qty{25}{\percent}.
Redirigir un usuario activado a un SDR captura por lo tanto un incremento relativo del
\qty{80}{\percent} en la probabilidad de conversión antes del primer contacto humano --- el
producto ya ha hecho la venta. La implicación operativa: el \emph{pipeline} de analítica de
producto debe hacer visible esta señal en el CRM en cuestión de minutos tras cruzar el umbral;
un retraso de 24 horas hace perder aproximadamente la mitad de la ventaja de conversión a medida
que la atención del usuario se desplaza a otros asuntos.
\end{example}

\begin{admonition}{Nota}
La decisión entre PLG y liderado por ventas no es una elección binaria a nivel de empresa. La
mayoría de las empresas por encima de \money{2000000} de ARR operan ambos movimientos en
paralelo, segmentados por tipo de comprador. Imponer un único movimiento a todos los segmentos
asigna mal los recursos: aplicar un SDR a un \emph{deal} de PYME de \money{600}/año destruye el
margen; dejar un prospecto empresarial de \money{50000}/año en una prueba de autoservicio hace
perder el \emph{deal} ante un competidor que tiene un defensor humano dentro de la cuenta. El
diagnóstico correcto son los indicadores unitarios de cada movimiento por segmento, no una
postura filosófica sobre PLG.
\end{admonition}

El movimiento GTM alimenta directamente la mecánica del funnel en \cref{ch:digital-funnel-and-crm}:
un movimiento PLG llena la parte superior del funnel con activaciones de prueba medidas por
\cref{eq:viral-k}; un movimiento liderado por ventas lo llena con transferencias de MQL desde
\emph{marketing}; PLS usa ambos. La valoración del \emph{pipeline} en \cref{eq:pipeline-ev}
aplica aguas abajo independientemente del movimiento, porque cada \emph{deal} en el CRM tiene un
valor nominal y una probabilidad de etapa tanto si llegó por autoservicio como por prospección del SDR.


\section{Gestión y compensación de la fuerza de ventas}\label{sec:salesforce-management}
% complejidad: 3

Construir un equipo de ventas requiere dos modelos separados: un \textit{modelo de
dimensionamiento} que traduce un objetivo de ingresos en un requisito de dotación de personal, y
un \textit{modelo de compensación} que alinea los incentivos del representante con los objetivos
de creación de valor de la empresa. Equivocarse en cualquiera de los dos es costoso. Un equipo
subdimensionado deja ARR sobre la mesa; un equipo sobredimensionado infla el gasto en ventas y
\emph{marketing} sin ingresos proporcionales. Un plan de comisiones mal diseñado recompensa el
comportamiento equivocado y, a través de la trampa de Kerr (\cref{sec:kerr-trap} en
\cref{ch:organizational-behavior}), produce sistemáticamente resultados no deseados por la empresa.

La fórmula de dotación de personal se deriva directamente del objetivo de ARR, la cuota por
ejecutivo de cuenta y la tasa histórica de cumplimiento:
\begin{equation}\label{eq:sales-productivity}
  \text{reps} \;=\;
    \left\lceil\frac{\text{ARR target}}{\text{quota}\times\text{attainment}}\right\rceil.
\end{equation}
La función techo garantiza que la empresa tenga suficiente capacidad para alcanzar el objetivo
incluso cuando los representantes no llegan al \qty{100}{\percent} de cumplimiento. El factor de
cumplimiento es el dato crítico: un plan construido sobre un \qty{100}{\percent} de cumplimiento
fallará cada trimestre; la tasa combinada histórica, típicamente entre \qty{70}{\percent} y
\qty{85}{\percent} para un equipo SaaS maduro, es el denominador correcto.

La compensación se estructura como ingresos objetivo totales (OTE) divididos entre un salario
base fijo y una parte variable vinculada al rendimiento:
\begin{equation}\label{eq:ote-structure}
  \text{OTE} \;=\; \text{base} \;+\; \text{variable}, \qquad
  \text{variable} \;=\; \text{quota} \times \text{commission rate}.
\end{equation}
Una distribución \qty{50}{\percent}/\qty{50}{\percent} entre base y variable es estándar para un
rol de AE empresarial; los roles de ventas internas y de SDR se acercan más a
\qty{70}{\percent}/\qty{30}{\percent} porque la influencia individual sobre el tamaño del
\emph{deal} es menor. La tasa de comisión en \cref{eq:ote-structure} es simplemente el objetivo
variable dividido entre la cuota, y establece la pendiente de la curva de incentivos.

\begin{example}[Dimensionamiento y costo del equipo de ventas para pasar de \money{4000000} a \money{6000000} de ARR]
La empresa tiene como objetivo \money{6000000} de ARR frente a una base actual de \money{4000000},
lo que requiere \money{2000000} en nuevo ARR. Cada ejecutivo de cuenta tiene una cuota de
\money{500000}; con la tasa histórica de cumplimiento del \qty{80}{\percent} combinado,
\cref{eq:sales-productivity} da:
\[
  \text{reps} \;=\;
    \left\lceil\frac{\money{2000000}}{\money{500000}\times0.80}\right\rceil
    \;=\; \lceil 5.00 \rceil \;=\; 5\ \text{AEs}.
\]
Cada AE tiene un OTE de \money{120000} con una distribución \qty{50}{\percent}/\qty{50}{\percent}:
base \money{60000}, variable \money{60000}. De \cref{eq:ote-structure}, la tasa de comisión es
\(\money{60000}/\money{500000} = \qty{12}{\percent}\) del ARR contratado. Cinco AEs representan
\money{600000} en compensación total --- la línea incremental de ventas y \emph{marketing} para
esta franja de crecimiento.

¿Supera la inversión el umbral LTV:CAC? Con un cumplimiento del \qty{80}{\percent}, cinco AEs
generan \(5\times\money{500000}\times0.80=\money{2000000}\) en nuevo ARR, es decir, aproximadamente
\num{3333} nuevos clientes al ACV de \money{600}/año de la hoja de datos. El CAC blended asistido
por ventas es \(\money{600000}/\num{3333}\approx\money{180}\). Frente al CLV de \money{3150}
proveniente de \cref{ch:customer-economics}, la razón LTV:CAC es
\(\money{3150}/\money{180}\approx17.5\times\), cómodamente por encima del piso de \(3\times\) de
\cref{eq:ltv-cac}. La restricción vinculante no es, por lo tanto, la economía, sino la ejecución:
encontrar cinco AEs que alcancen su cuota.
\end{example}

\begin{admonition}{Advertencia}
Las cláusulas de aceleración que multiplican las comisiones por encima del \qty{100}{\percent} de
cumplimiento de cuota crean un incentivo racional para que los representantes guarden
oportunidades al inicio del trimestre --- acumulando \emph{deal}s hasta tener suficiente
\emph{pipeline} para garantizar el alcance del umbral del acelerador antes de presentarlos. La
empresa entonces observa patrones de contratación irregulares concentrados al final del trimestre
que comprimen la cola de implementación e incorporación, elevan el \emph{churn} en el trimestre
siguiente y hacen el pronóstico poco fiable. La solución estructural es pagar la comisión sobre
el ARR \textit{retenido} al final del primer período de renovación, en lugar de hacerlo sobre el
ARR contratado en el momento de la firma: un representante cuyo cliente hace \emph{churn} en los
primeros doce meses devuelve una parte de la comisión inicial. Esto aborda directamente la trampa
de Kerr (\cref{sec:kerr-trap}): pagar sobre ARR contratado recompensa el cierre mientras se
espera que la retención en la renovación llegue sola. \textcite{dixon2011challenger} documenta la
taxonomía completa de distorsiones en el comportamiento de ventas que emergen de un diseño de
comisiones mal alineado.
\end{admonition}

La fórmula de productividad de ventas y la estructura OTE son datos de entrada al modelo operativo
en \cref{ch:operating-model}, donde aparecen como partidas en la proyección de gastos de ventas y
\emph{marketing}. Vincular el diseño de la compensación al \emph{churn} garantiza que el
crecimiento de ARR modelado en \cref{eq:mrr-growth} no sea erosionado por el \emph{churn}
temprano inducido por ventas.


\section{Ventas B2B y gestión del \emph{pipeline}}\label{sec:b2b-selling}
% complejidad: 3

Un equipo de ventas que no puede pronosticar sus propios ingresos es un pasivo, no un activo. El
\emph{pipeline} en bruto --- la suma de los valores nominales de todos los \emph{deal}s abiertos
--- es la métrica de vanidad canónica de las ventas empresariales: equipara un \emph{deal} en la
primera llamada de descubrimiento con uno en revisión legal final. El \emph{pipeline} ponderado
reemplaza el valor nominal por el valor esperado de cada oportunidad, donde la expectativa está
condicionada a la etapa del funnel en que se encuentra el \emph{deal}:
\begin{equation}\label{eq:pipeline-forecast}
  \widehat{\text{ARR}} \;=\; \sum_{i=1}^{n} v_i \cdot p_{\,\text{stage}(i)},
\end{equation}
donde \(v_i\) es el valor nominal del \emph{deal} \(i\) y \(p_{\,\text{stage}(i)}\) es la
probabilidad histórica de cierre para su etapa actual. \Cref{eq:pipeline-forecast} es el análogo
en ARR de \cref{eq:pipeline-ev} en \cref{ch:digital-funnel-and-crm}: ambos aplican el operador
de valor esperado a una cartera de conversiones probabilísticas. Las probabilidades de etapa deben
calibrarse con las tasas de victorias observadas en cada etapa, no asignarse por optimismo
directivo; un modelo de \emph{pipeline} construido sobre probabilidades de etapa infladas falla a
la misma tasa que no tener modelo, con el daño adicional de generar una falsa confianza.

Junto con la instantánea de valor esperado, la Ley de Little de \cref{eq:littles-law} proporciona
la duración media del ciclo de ventas a partir del inventario de \emph{pipeline} y el rendimiento:
\[
  W \;=\; \frac{L}{\lambda},
\]
donde \(L\) es el número de \emph{deal}s abiertos, \(\lambda\) es el número de \emph{deal}s
cerrados por período y \(W\) es el tiempo medio que un \emph{deal} permanece en el
\emph{pipeline}. Un \(W\) creciente es una señal de alerta temprana de que los \emph{deal}s se
están estancando, no de que el \emph{pipeline} sea saludable.

La secuencia de preguntas SPIN (\textcite{rackham1988spin}) y la metodología Challenger
(\textcite{dixon2011challenger}) operan a nivel de la llamada de ventas individual, pero producen
la progresión de etapas que alimenta \cref{eq:pipeline-forecast}. El arco situación-problema-implicación-beneficio
del SPIN hace aflorar el dolor económico que justifica una decisión de compra; el arco
enseñar-adaptar-tomar-control del Challenger reencuadra el modelo mental del comprador antes de
que un competidor pueda anclar su posición. Ambos marcos existen para mover los \emph{deal}s hacia
la derecha a través de las etapas del \emph{pipeline}, incrementando \(p_{\,\text{stage}(i)}\)
para la cartera en curso.

\begin{example}[Pipeline ponderado y Ley de Little]
El equipo de ventas tiene \num{40} \emph{deal}s abiertos distribuidos uniformemente en cuatro
etapas del \emph{pipeline}: \num{10} en Descubrimiento, \num{10} en Demostración, \num{10} en
Propuesta y \num{10} en Negociación. Cada \emph{deal} tiene un valor nominal de \money{8250} de
ARR. Las tasas históricas de cierre por etapa son \qty{15}{\percent} en Descubrimiento,
\qty{35}{\percent} en Demostración, \qty{60}{\percent} en Propuesta y \qty{85}{\percent} en
Negociación. Aplicando \cref{eq:pipeline-forecast}:
\begin{align*}
  \widehat{\text{ARR}}
    &= 10\times\money{8250}\times0.15
    \;+\; 10\times\money{8250}\times0.35 \\
    &\quad +\; 10\times\money{8250}\times0.60
    \;+\; 10\times\money{8250}\times0.85 \\
    &= \money{12375} + \money{28875} + \money{49500} + \money{70125}
     \;\approx\; \money{161000}.
\end{align*}
El valor nominal bruto del \emph{pipeline} es \(40\times\money{8250}=\money{330000}\) --- más del
doble del ARR esperado de \money{161000}. Un gerente de ventas que reportara el número bruto como
«cobertura de \emph{pipeline}» estaría sobreestimando el pronóstico real en un \qty{105}{\percent}.

Aplicando la Ley de Little: si el equipo cierra \num{10} \emph{deal}s por mes, el \emph{deal}
promedio permanece \(W = 40/10 = 4\) meses en el \emph{pipeline}. Si el equipo añade \num{5}
oportunidades generadas por SDR por mes pero cierra al mismo ritmo de \num{10}/mes, \(L\) crece a
\num{60} dentro de cuatro meses y \(W\) se extiende a \(60/10 = 6\) meses. El \emph{pipeline}
parece más saludable (más \emph{deal}s), pero la fecha esperada de cierre de cualquier
\emph{deal} dado se ha demorado dos meses --- un problema de pronóstico que se disfraza de éxito
del \emph{pipeline}.
\end{example}

\begin{admonition}{Nota}
Los marcos de calificación --- como MEDDPICC (Métricas, Comprador Económico, Criterios de
Decisión, Proceso de Decisión, Proceso Documental, Identificación del Dolor, Defensor,
Competencia) --- cumplen para el \emph{pipeline} la misma función que la puntuación de segmentos
(\cref{eq:segment-score}) cumple para el mercado direccionable: obligan a realizar una estimación
explícita y medible de la calidad de la oportunidad antes de comprometer capacidad de ventas. Una
oportunidad que no supera la lista de verificación de calificación debe retroceder de etapa o
descalificarse, en lugar de mantenerse con una probabilidad optimista; el costo es un
\emph{pipeline} temporalmente más pequeño, y el beneficio es un \cref{eq:pipeline-forecast} en
el que la dirección puede confiar realmente.
\end{admonition}

La estructura de canales, el movimiento GTM, el diseño de la fuerza de ventas y la gestión del
\emph{pipeline} son las cuatro capas de decisión que convierten un producto y un segmento objetivo
en ingresos predecibles. Cada capa alimenta a la siguiente: el canal determina la estructura de
costos de adquisición (\cref{eq:channel-margin}), el movimiento determina si ese costo está
dominado por el costo fijo de ventas o por la inversión variable en producto, el diseño de la
compensación determina si el comportamiento de los representantes se alinea con la maximización
del LTV (\cref{eq:ltv-cac}), y el modelo de \emph{pipeline} determina si la organización puede
leer su propio progreso. Las cuatro capas alimentan el funnel digital en
\cref{ch:digital-funnel-and-crm} y, en última instancia, la línea de ingresos del modelo
operativo en \cref{ch:operating-model}.
